\documentclass[10pt,a4paper]{amsart}
\usepackage[margin=19mm]{geometry}
\usepackage{amsmath,amssymb,mathtools}
\usepackage[scr=rsfs,cal=euler]{mathalfa}
\usepackage{xfrac}
\usepackage[unicode,hidelinks,bookmarks=false]{hyperref}
\newcommand{\ie}{i.e.\ }
\newcommand{\cf}{cf.\ }
\newcommand{\resp}{respectively\ }
\newcommand{\Subsection}{\subsection}
\providecommand{\nobreakdash}{\nobreak\hbox{-}}
\setlength{\emergencystretch}{2em}
\multlinegap0pt
\usepackage{bm}
\usepackage{bbm}
\usepackage{dsfont}
\usepackage{xspace}
\usepackage{cancel}
\usepackage{mathtools}
\usepackage{amsthm}
\makeatletter
\@ifundefined{theorem}{\newtheorem{theorem}{Theorem}}{}
\makeatother
\def\Sc{\ensuremath{\mathscr{S}}}
\def\bK{\ensuremath{\mathbb{K}}}
\def\bR{\ensuremath{\mathbb{R}}}
\def\cal{\ensuremath{\boldsymbol{\circlearrowleft}}}
\title[Explicit Families of Spinor Representations]{Explicit Families of Spinor Representations}
\author{Jesus Sanchez Jr}
\date{Source: arXiv:2505.16203v1 (CC BY 4.0)}
\begin{document}
\maketitle

\noindent\emph{Source and licence.} Explicit Families of Spinor Representations. Version arXiv:2505.16203v1 (CC BY 4.0), distributed under the Creative Commons Attribution 4.0 International licence, as recorded in the arXiv licence metadata for this exact version and shown on the abstract page https://arxiv.org/abs/2505.16203v1. Full rights notice: licenses/arxiv-2505.16203-CC-BY-4.0.txt.\par\smallskip

\noindent\textit{Editorial scope.} Counted unit: the Theorem whose source label is \texttt{None} in the article (source0.tex lines 634--636, source label \texttt{None}), delivered together with the article's own complete original proof of it (source0.tex lines 637--656, one \texttt{proof} environment of 20 lines). No mathematics is written, repaired or re-derived: statement and proof are the article's own text line for line. Required context retained uncounted: none. Every definition, display and label of those lines is retained unchanged. Omitted from this package and not redistributed: 1--633, 657--1230. Nothing inside a retained excerpt is omitted or altered: every retained body is delivered exactly as the article's lines are written, so no omission receipt of any kind accompanies this package. Citations: the retained excerpts carry no citation command at all, so no bibliography entry of the article is transcribed and no reference is redistributed. Reference adaptation: none was needed and none was made; every reference inside the delivered unit resolves inside it, so no unresolved reference or citation is produced. Printed slips kept literal: no printed word, symbol, hypothesis or number of the counted statement or of its proof was altered. Layout and class: the article's own class is \texttt{article} with packages including inputenc, amsmath, amscd, amssymb, amsfonts, amsthm, courier, relsize, bm, yhmath, hyperref, enumerate, mathrsfs, mathtools; the wrapper is the lane's pinned \texttt{amsart} editorial wrapper at 10pt on A4, which restates the theorem-like environments the retained text needs and adds the article's own macro definitions that the retained text needs (\texttt{\textbackslash Sc}, \texttt{\textbackslash bK}, \texttt{\textbackslash bR}, \texttt{\textbackslash cal}), with extra wrapper packages (bm, bbm, dsfont, xspace, cancel, mathtools). The delivered numbering is the unit's own. Assets: no retained span contains a figure environment, a graphics-inclusion command, a PDF-page inclusion, a table or a TikZ picture; no image file is copied into this package, the package declares no asset dependency and no image file is redistributed with it. Licence: version arXiv:2505.16203v1 of this article is distributed under the Creative Commons Attribution 4.0 International licence, as recorded in the arXiv licence metadata for this exact version and shown on the abstract page https://arxiv.org/abs/2505.16203v1. A full rights notice is delivered at licenses/arxiv-2505.16203-CC-BY-4.0.txt. Characters: every retained excerpt is pure ASCII, so the delivered engine, XeLaTeX with its default Latin Modern text font, reports no missing character.
\begin{theorem}
    An oriented Euclidean vector bundle $(E,g_E,\cal_{\bR})\rightarrow M$ of rank $n$ admits a spin structure if and only if there is a real Euclidean vector bundle $(\Sc_E,g_{\Sc})\rightarrow M$ which is a bundle of $C\ell(E,g_E)-\underline{\bK_n}$ bimodules as well as a bundle of $C\ell^0(E,g_E)-\underline{\bK^0_n}$-bimodules for which the Clifford action is irreducible and the Clifford action by vectors in $E$ and right multiplication by $a\in\emph{Im}\underline{\bK_n}$ is skew-adjoint.
\end{theorem}

\begin{proof}
    Suppose that $(E,g_E,\cal_{\bR})\rightarrow M$ admits a spin lift $P\circlearrowleft Spin(n)\rightarrow P_{SO}(E,g_E,\cal_{\bR})\circlearrowleft SO(n)$. If $S_n$ denotes the real irreducible module constructed for $C\ell_n$ above, note that the spin representation $\rho :Spin(n)\rightarrow GL(S_n,\bR)$ is given by restriction of the action $C\ell_n\circlearrowright S_n$ to $Spin(n)$. In particular, the spin representation inherits all of the extra bells and whistles that the Clifford action has, i.e. commutativity with the right action by $\bK_n$ and the real metric $g_{S_n}$ for which the left action by vectors in $\bR^n$ and left action by elements in $\text{Im}\bK_n$ is skew-adjoint. Moreover, note that $Spin(n)\subseteq C\ell^0_n$, hence if $C\ell^0_n$ has a larger right intertwiner algebra $\bK^0_n$, then the action $Spin(n)\circlearrowright S_n$ will be compatible with the right action by $\bK^0_n$ as well. Since the right action $C\ell_n\circlearrowright S_n$ is compatible with the respective spin representations, this action descends to an action
    \[
        P\times_{Ad\rho} C\ell_n\circlearrowright (P\times_{\rho}S_n,g_{S_n})\circlearrowleft \underline{\bK_n},\quad P\times_{Ad\rho}C\ell^0_n\circlearrowright P\times_{\rho}S_n\circlearrowleft \underline{\bK^0_n}
    \]
    Because $C\ell(E,g_E)\simeq P\times_{Ad\rho}C\ell_n$ as a bundle of algebras, we have the asserted spinor module.\par
    Conversely, suppose we are given such a spinor module 
    \[
        C\ell(E,g_E)\circlearrowright (\Sc_E,g_{\Sc})\circlearrowleft\underline{\bK_n},\quad C\ell^0(E,g_E)\circlearrowright \Sc_E\circlearrowright \underline{\bK^0_n}
    \]
    We define the \textbf{\emph{principal bundle of spin coordinate systems for $(E,g_E,\cal_{\bR},\Sc_E)$}} to be the disjiont union
    \[
    P_{Spin}(E,g_E,\cal_{\bR},\Sc_E):=\bigsqcup_{m\in M}P_{Spin}(E_m,g_{E_m},\cal_{\bR},\Sc_{E_m}).
    \]
    The smooth structure is the natural choice and this fibration $P_{Spin}(E,g_E,\cal_{\bR},\Sc_E)\rightarrow M$ admits a right free and fibered transitive action of $Spin(n)$ such that the natural map
    \[
        p:P_{Spin}(E,g_E,\cal_{\bR},\Sc_E)\rightarrow P_{SO}(E,g_E,\cal_{\bR})
    \]
    is a smooth, fibered double covering which is compatible with the double covering $Ad:Spin(n)\rightarrow SO(n)$. Thus, we have obtained the desired spin structure $P_{Spin}(E,g_E,\cal_{\bR},\Sc_E)$.
\end{proof}

\end{document}
